\begin{Solbook}{1}
 \par \noindent  a) $x=1\,\text {và}\,x=2$ \hspace {2cm} b) $x<0$  \hfill \qedEX 
\end{Solbook}
\begin{Solbook}{2}
 \par \noindent Nếu làm một mình thì đội 1 hoàn thành công việc sau 20 ngày, đội 2 hoàn thành công việc sau 30 ngày. \hfill \qedEX 
\end{Solbook}
\begin{Solbook}{3}
 \par \noindent  Chiều dài $AB$ của mặt cầu thang khoảng $6$ m.  \hfill \qedEX 
\end{Solbook}
\begin{Solbook}{4}
 \par \noindent  a) $A=6\sqrt {3}$ \hspace {2cm} b) $B=5$ \hspace {2cm} c) $C=-3\sqrt {x}$  \hfill \qedEX 
\end{Solbook}
\begin{Solbook}{5}
 \par \noindent  \immini { \begin {enumerate} \item $\widehat {AOB}=120^\circ ,\widehat {AKB}=120^\circ $. \item Chứng minh tam giác $MAB$ đều và $K$ là tâm đường tròn ngoại tiếp tam giác $MAB$. \item Tam giác $MAB$ đều nên $AK$ vuông góc với $BM$. Hơn nữa $OM$ vuông góc với $AB$. \end {enumerate} }{ \begin {tikzpicture}[line width=0.8pt, scale=1] \def \R {2.5} \def \A {5} \path (0,0) coordinate (O) (\A ,0) coordinate (M) ($(O)!0.5!(M)$) coordinate (K); \par \path [name path=ciro] (O) circle (\R ); \path [name path=cirom] (K) circle (\A /2); \par \path [name intersections={of=ciro and cirom, by={A,B}}]; \par \draw (O) circle (\R ) (A)--(O)--(B) (O)--(M) (M)--(B)--(A)--cycle; \par \path [name path=om] (O)--(M); \coordinate (H) at (intersection of A--B and O--M); \coordinate (At) at ($(A)!2!(K)$); \coordinate (I) at (intersection of A--At and B--M); \coordinate (Y) at ($(A)!0.5!(M)$); \draw (Y) let \p 2=($(Y)-(A)$) in circle ({veclen(\x 2,\y 2)}); \draw (A)--(K)--(B) (K)--(I)--(H); \foreach \x /\g in {O/180,M/-40,A/90,B/-60,K/45,I/-90,H/-135} \draw [fill=white] (\x ) circle (1.5pt) +(\g :.3) node{$\x $}; \end {tikzpicture} }  \hfill \qedEX 
\end{Solbook}
\begin{Solbook}{6}
 \par \noindent  Biến đổi hệ phương trình ban đầu thành $\begin {cases} b^3-3a^2b=1\\ a^3-3ab^2=2 \end {cases}$ \par Bình phương hai vế hai phương trình của hệ \par Cộng hai vế của hai phương trình vừa tìm được ta được $(a^2+b^2)^3=5$ \par Từ đó suy ra $a^2+b^2=\sqrt [3]{5}$  \hfill \qedEX 
\end{Solbook}
